\chapter{}

Pietro, con gratitudine, sentiva vicino a Ghìsola le sue prime emozioni delicate. Ammirò un fiore quando gli venne voglia di coglierlo per lei; e, non arrischiandosi, lo buttava via; quando era ancora per non crederci, provando una diminuzione di sè stesso. E come tutta la natura gli apparve a un tratto misteriosa, con violenza! Qualche cosa da disperarsene!

Era stato bocconi in terra, chiudendo tra le braccia un pulcino per tenerlo con sè! Aveva aiutato le formiche, togliendo dal loro cammino un bastone che dovevano valicare esitando e poi disperate: tremolando con un chicco troppo grosso, che le faceva cadere capovolte! Teneva con tenerezza un indovinello in mano, e lo rimproverava quando volava via!

Cercava di superare le sue malinconie; ma non poteva dimenticarle quanto avrebbe voluto. Talvolta ne era distaccato di soprassalto; e allora gli veniva uno stato mentale confuso e torbido che pareva sempre per andarsene. E aveva l’illusione che il suo spirito assumesse così enormi proporzioni che i suoi pensieri vi si smarrivano dentro, insieme con i loro echi improvvisi, come in una stanza troppo grande. Quante volte non s’era considerato perduto, mentre le imagini esteriori lo invadevano senza tregua! Ora gli pareva di avere la propria anima; ora diminuiva; mentre questi movimenti gli davano un malessere come quello delle vertigini.

Talvolta gli pareva di trovarsi a scuola dove tutto a un tratto entrava una grancassa; e allora si sentiva tanta voglia di ridere che si spaventava, soffocando il grido dell’incubo. Anna credeva che avesse male; e gli metteva una mano su la fronte, dicendogli:

— Ti viene la febbre?

Egli gridava, allora:

— No! No! Lasciami stare!
